\section{模型求解、检验与敏感性分析}

\subsection{求解流程与最优性}

程序使用Python读取附件，调用HiGHS分支定界算法逐日求解问题一、问题二，再求解月度0--1排班。日模型规模较小，可直接要求相对最优间隙为0；月度模型在理论下界581人处即获得可行解，因此结合式\eqref{eq:lb}可证明问题三全局最优。员工排班是经典组合优化问题，滚动窗口建模可直接表达连续工作限制\cite{ernst2004,pinedo2022}。

\subsection{独立可行性检验}

求解后不依赖求解器状态，程序重新执行以下检查：
\begin{enumerate}[label=\arabic*)]
  \item 每天恰有5个非空班次，班次人数为非负整数；
  \item 每小时实际处理量不超过相应班次提供的能力；
  \item 逐小时重新计算$b_{dh}=\sum_{\tau=0}^{h}(a_{d\tau}-p_{d\tau})$，确认$b_{dh}\ge0$且$b_{d,23}=0$；
  \item 问题二确认每名工人恰有1个低效率小时，并验证式\eqref{eq:deadline}；
  \item 问题三逐人核对工作天数与全部23个连续8天窗口，并逐日核对覆盖人数。
\end{enumerate}
全部检查通过。详细变量值写入\texttt{results.json}，班次和人员编号写入\texttt{排班结果.xlsx}。

\subsection{处理效率敏感性}

将正常和低效率处理能力同时上下扰动5\%，重新求解问题二全部30天，结果见表\ref{tab:sensitivity}。

\begin{table}[htbp]
\centering
\caption{问题二对处理效率的敏感性}
\label{tab:sensitivity}
\begin{tabular}{crrr}
\toprule
效率系数 & 总人日 & 峰值人数 & 相对基准变化\\
\midrule
0.95 & 12576 & 611 & $+5.23\%$\\
1.00 & 11951 & 581 & --\\
1.05 & 11383 & 553 & $-4.75\%$\\
\bottomrule
\end{tabular}
\end{table}

效率下降5\%使总人日增加5.23\%、峰值增加30人，略高于反比例关系，原因是人数整数化、班次覆盖和16时截止约束共同放大了局部高峰。由此建议实际部署保留约5\%机动能力，或在预测到早间高峰时提前调整班次起点。
